\section{A compact controller and its actual deployment envelope}\label{app:unified}
It is possible to keep every installed parameter fixed through an entire run without enumerating every observed transcript. This requires a particular structure of the conditional source rows. The four results below give the representation, the approximation certificate, a source satisfying the linked physical premises, and its limits.

\begin{theorem}[U.G1. Sufficient counts without a world prior]\label{adv:ug1}
Fix horizon \(H\), finite alphabets of commands \(A\), observations \(O\), and public states \(S\). A world fixes one tag \(m\) and a continuous parameter \(\theta\) in that tag's compact rational semialgebraic domain. For tag \(m\) the hidden start is fixed and hidden states form a finite set \(Z_m\). The next public state is a deterministic function of time, current public state, action and received letter alone. Any separately revealed public update is included in that letter and its likelihood. The next hidden state for each tag is deterministic given the same inputs and its current hidden state, independent of \(\theta\). The letter probability is a specified member of \(r_m\) polynomial factors \(f_{mj}(\theta)\); rows are normalized and nonnegative. Include complete relevant responder memory and flags. Let \(q\) retain time, public state and each tag's propagated hidden state and factor counts \(n_{mj}\). Write \(L_m(q,\theta)=\prod_j f_{mj}(\theta)^{n_{mj}}\), using \(0^0=1\) for an unused factor. Keep a letter whenever some actual world's prefix likelihood through it is positive.

Require that the common callable actions, stops, supported-path constraints, all relevant observations and the conditional expectations of each additive score are determined by \(q\) in every surviving world. Then nonnegative flows \(F(q),F(q,a)\) give exactly the attainable simultaneous expected-performance vectors. Root inflow is one, outgoing action flows sum to node inflow, and each incoming letter edge from \((p,a)\) contributes \(F(p,a)\) once, even for letters sharing a destination. Multiply the appropriate inflow or action flow by the world likelihood and conditional score when evaluating that world. The graph has at most
\[
 (H+1)|S|\prod_m\left(|Z_m|\binom{H+r_m}{r_m}\right)
\]
nodes and \(|A||O|\) edges per node.
\end{theorem}
\begin{proof}
Fix a complete history \(h\) and call \(x(h)\) the private-seed mass making its successive decisions if presented with those data. In world \((m,\theta)\) the history mass factors as \(x(h)L_m(q,\theta)\). This is where source nonanticipativity and independence of private seed are used. Sum these seed masses, and the masses assigned to each next action, over histories with the same \(q\). Each distinct possible next letter carries the \emph{same} action-seed mass into its corresponding successor. The chance of that letter is already in the new likelihood. Conversely choose an action with probability \(F(q,a)/F(q)\) at positive inflow. Layer induction recovers the flows. Zero inflow requires no ratio and no physical move. A totalized implementation supplies a lawful default only where its interface requires and provides one. Summing likelihood-weighted conditional scores recovers every worldwise performance simultaneously, with hard validity tested on their whole union of supports. Flows at depth \(t\) sum to at most \(|O|^t\), hence the closed feasible flow polytope is bounded. Counting factor multisets of size at most \(H\) yields the graph bound. Factor degree \(d\) gives likelihood degree at most \(Hd\) before multiplying the score. Exact support checking, arithmetic precision and continuous robust optimization are additional costs.
\end{proof}

\begin{theorem}[U.G2. A lower and upper deployment value]\label{adv:ug2}
Let \(X\) be a nonempty compact rational realization polytope, exact for both implementations including all hard constraints and exhaustive for the policies compared. Partition actual worlds into finitely many labeled compact strata. In stratum \(j\), use a finite net of actual worlds with verified radius \(r_j\). For an actual affine constraint \(g_i^D(x,w)\le b_i\), suppose its reference version \(g_i^R\) has uniform Lipschitz constant \(L_{ij}\) in \(w\) and its discrepancy from actual is at most \(d_{ij}\). Make the analogous assumptions for bounded cost \(c\), with constants \(L_{cj},d_{cj}\). Write \(\alpha_{ij}=L_{ij}r_j,\ \alpha_{cj}=L_{cj}r_j\). Use exact rational net coefficients or certified one-sided objective bounds.

At net points relax the constraint to \(g_i^R\le b_i+d_{ij}\) and use cost \(c^R-d_{cj}\) for a lower minimum \(L\). Tighten to \(g_i^R\le b_i-d_{ij}-\alpha_{ij}\) and use \(c^R+d_{cj}+\alpha_{cj}\) for an upper minimum \(U\). Minimize the worst net cost in each case. Lower infeasibility proves actual infeasibility. Any feasible upper solution is a physical policy with cost at most its reported upper objective. If both are feasible, \(L\le C^*\le U\), where \(C^*\) is the actual minimax optimum. If one controller has uniform actual constraint margins \(\gamma_{ij}>0\) and \(t=\max_{ij}(2d_{ij}+\alpha_{ij})/\gamma_{ij}\le1\), then
\[
 U-L\le tR+\max_j(2d_{cj}+\alpha_{cj})
\]
for reference costs in a common interval of width \(R\).
\end{theorem}
\begin{proof}
An actual feasible policy passes every relaxed net test and its actual worst cost bounds the lower objective. From any tightened solution move to the actual world through its reference net point, spending at most \(d_{ij}+\alpha_{ij}\) on a constraint and the analogous amount on cost. For the gap bound, mix fraction \(t\) of the strictly feasible policy into a lower optimizer. At a net point the mixture scores at most \(b_i+d_{ij}-t\gamma_{ij}\), which meets the tightened cap. Cost changes by at most \(tR\) before the two objective allowances. A common sublaw of mass \(1-d\) preserving true issuance and violation changes the signed selective score by at most \(d\) because its range has width one. It cannot transfer hard authority. An empty upper problem can reflect zero slack in a feasible actual problem.
\end{proof}

\begin{theorem}[U.G3. The guarded common family]\label{adv:ug3}
Use the source of Theorem~\ref{body:unified}. The installed state is \((\theta,m,k,B,\beta)\), with the two \(\theta\) intervals specified there, hidden bit \(m\), route \(k\in\{L,R\}\), guard revision \(B\), and its allowed surcharge \(\beta\). Starting at \(v=0\), each of \(N\) received prelude bits stays at the current \(v\) with conditional probability
\[
 p_m(\theta)=\tfrac12+(-1)^m\theta(1-\theta)/2,
\]
otherwise it flips; the new bit becomes \(v\). Actual safe intake earns \(1/2\) utility and reports the only callable route. The optional probes of Section~\ref{sec:unified} use fresh bits, with no repeat or wrong-route action. After a positive report the carried responder proposes \(a=m\mathbin{\mathrm{xor}}v\), then the finalizer draws independently at the threshold \(\theta/4+a/16+\beta\). A true positive safe successful routed finish adds \(1/2\) utility. A negative report stops. Other admitted startup and disclosure channels give no extra separation on the least favorable pair. Exact source authority and event mappings cover every reached cut. Guard check, comprehensive inspection and ordinary-route read respectively cost \(1/4,1,1/8\); only the first two separate the actual guard revisions.

On a good receipt Theorem~\ref{body:unified}'s optimum and bounds hold, including nonzero true risk. For a common pre-declaration process with a nominal zero surcharge, total variation of the compared decoded laws is at most \(\beta\Pr(I)\le1/32\), and the increase in violation mass is exactly \(\beta\Pr(I)\). At depth \(t\) there are \(t+1\) sufficient prelude states against \(2^t\) supported words. The graph, the grid bounds and the optional calibration conclusions use the distinct source hypotheses below.
\end{theorem}
\begin{proof}
The matched endpoints \(\theta=1/4,3/4\) have the same \(\theta(1-\theta)\), so every prelude, receipt, clock and other admitted pre-probe field can be coupled identically at matched \(m,k,B,\beta\). The proof of Theorem~\ref{body:unified} now applies to all lawful private choices. Its mix issues with probability \(q=5/6-2\theta/3\). The restricted high-or-skip optimum buys high with probability \(2/3\) and guesses fairly otherwise, giving the same \(q\). Every good issued cut has violation probability between \(1/32\) and \(5/16\), with upper \(5/32\) for positives. The proposal changes the physical conditional probability by \(1/16\). Actual intake and safe positive finishing give the utility calculation in Section~\ref{sec:unified}. A true-positive bad world in any unresolved guard cell must receive service with positive probability and has conditional risk at least \(1/2+1/32\); the ordinary check cannot rule it out. This proves the least realized setup price within the given catalog.

Couple nominal and actual to share all earlier randomness, lawful commands and the proposal. At an issued cut the single uniform draw lies between the two thresholds with probability \(\beta\). Only those outcomes disagree on the declared score fields. This gives both the total-variation bound and the exact increment. Any included upstream calibration must have the same actual law on both sides, or the comparison must condition on its common carried record or include its own discrepancy. Uncompared raw fields require new evidence.

At depth \(t\), the number of stays determines the number of flips, final bit and both tags' likelihood products. Every word has positive probability; retain true authority and the recorded route in the later constant number of stages. Theorem~\ref{adv:ug1} thus gives a quadratic total node count. Within a fixed stratum a change \(\Delta\theta\) changes each prelude row by at most \(|\Delta\theta|/2\), the probe by at most \(|\Delta\theta|\), and finish by at most \(|\Delta\theta|/4\). First-discrepancy coupling gives Lipschitz constant \(N/2+5/4\) for width-one scores. When the prefix is independent of the surcharge, issuance is constant in \(\beta\). The signed risk and negative utility can then be checked at \(\beta=1/32\). For this exact endpoint evaluation an always-high policy has the requisite margin on the nonautomatic constraints: \(1/12\) for error and positive coverage, \(1/128\) for signed risk, and \(9/256\) for positive utility. A mesh radius no larger than \(1/[128(N/2+5/4)]\) gives a feasible upper controller. Refine under Theorem~\ref{adv:ug2} for a shrinking interval. The symbolic cost proof uses no grid.

If an authorized optional calibration returns bounded Bernoulli observations of conditional mean \(\beta\), Theorem~\ref{adv:g023} with one stream, zero bias and separately reserved cost gives a simultaneous confidence bound on \(\beta\). Report a sharper margin only if it holds for every parameter in the intersection of the confidence set and verified guard envelope. Its correctness then follows on coverage. A prefix whose law depends on \(\beta\) may change an arbitrary controller's issuance, so the endpoint projection above does not extend automatically to it. Count such a prefix in the matched price only with the same capped interface and matched opposite laws on both sides. Actual hard authorization still comes from the guard.
\end{proof}

\begin{theorem}[U.G4. Stopping and information limits]\label{adv:ug4}
For an almost surely terminating policy with integer stopping time \(\tau\), assume the uniform law \(\sup_w\Pr_w(\tau>n)\le b_n\), a lawful fallback at cutoff \(H\), zero additional fallback cost, and original nonnegative per-step cost at most \(K\). Truncation changes expected cost by at most \(K\sum_{n\ge H}b_n\), and a common width-\(D\) terminal metric by at most \(Db_H\). For a discounted per-step score bounded in absolute value by \(K\), zero future fallback contributions bound the error by \(K\lambda^H/(1-\lambda)\) for discount \(0<\lambda<1\).

Separately let \(m\) actual bits determine identity. If bit \(i\) has a lawful repeatable bounded unit-interval read of price \(c_i\), conditional mean always at most \(\ell_i\) or always at least \(\ell_i+\gamma_i\) according to that bit, a lawful schedule of \(n_i=\lceil2\gamma_i^{-2}\log(m/\delta)\rceil\) reads per bit identifies the contract with error at most \(\delta\) at price \(\sum_i c_i n_i\). Conversely, if for every \(\kappa>0\) some opposite pair has startup and all corresponding conditional next-record laws within total variation \(\kappa\) on a common lawful finite horizon, no audit at that horizon can uniformly beat chance in the minimax binary sense.
\end{theorem}
\begin{proof}
Share the path until termination or cutoff. The integer tail sum bounds \(\mathbb E[(\tau-H)_+]\), and a terminal discrepancy can occur only on the tail. Sum a geometric series for the discounted bound. Arbitrarily late deterministic stops across different worlds have no uniform tail estimate. Even in one world a tail \(1/(n+1)\) has infinite mean. Discount alone can obscure a late undiscounted violation. For the identity result, threshold each empirical mean at \(\ell_i+\gamma_i/2\). The conditional bounded-draw inequality in Theorem~\ref{adv:g023} bounds its one-sided error by \(\exp(-n_i\gamma_i^2/2)\); summing proves the claim. This upper bound excludes exponential identity price when the schedule, gaps and prices give a polynomial total. Finally couple startup and each reached letter while sharing the controller's private seed. Record discrepancy is at most \(\min(1,(H+1)\kappa)\). The two errors sum to at least one minus that distance. Let \(\kappa\) decrease to zero. This argument pertains to the complete lawful interface, including all initial observations.
\end{proof}
